\chapter{Esercizi Finanza}
\section{Esercizi VAN e VA}

\subsection{Richiamo delle formule}
\subsection*{Formule utilizzate}

\noindent\textbf{1. Valore attuale di più flussi di cassa}

\[
VA=\sum_{n=1}^{N}\frac{F_n}{(1+i)^n}
\]

\noindent dove:
\begin{itemize}
    \item $F_n$ = flusso di cassa ricevuto nell'anno $n$
    \item $i$ = tasso di attualizzazione
    \item $n$ = anno in cui si verifica il flusso
\end{itemize}

\noindent\textbf{2. Valore Attuale Netto (VAN)}

\[
VAN=F_0+VA
\]

\noindent oppure, in forma completa:

\[
VAN=F_0+\sum_{n=1}^{N}\frac{F_n}{(1+i)^n}
\]

\noindent dove $F_0$ rappresenta l'investimento iniziale ed è generalmente
negativo perché costituisce un'uscita di denaro.

\noindent\textbf{3. Valore attuale di una rendita annua costante}

\[
VA=C\left[
\frac{1}{r}
-\frac{1}{r(1+r)^t}
\right]
\]

\noindent dove:
\begin{itemize}
    \item $C$ = flusso di cassa costante ricevuto ogni periodo
    \item $r$ = tasso di attualizzazione
    \item $t$ = numero di periodi
\end{itemize}

\noindent Questa formula è una scorciatoia utilizzabile quando si riceve
\textbf{lo stesso flusso di cassa per un numero determinato di periodi}.
Una volta trovato il valore attuale della rendita, il VAN si calcola come:

\[
VAN=-\text{Investimento iniziale}+VA
\]
\newpage
\subsection{Esercizio 1}

La società \textbf{InnovaTech S.r.l.} sta considerando di investire in un progetto di espansione aziendale, caratterizzato dai seguenti flussi di cassa:

\begin{itemize}
    \item \textbf{Investimento iniziale:} 100.000 euro.
\end{itemize}

\noindent I flussi di cassa netti previsti per gli anni successivi sono:

\begin{itemize}
    \item Anno 1: 30.000 euro;
    \item Anno 2: 40.000 euro;
    \item Anno 3: 50.000 euro;
    \item Anno 4: 60.000 euro;
    \item Anno 5: 70.000 euro.
\end{itemize}

\noindent Il tasso di sconto appropriato per questo tipo di investimento è del \textbf{10\% annuo}.

\medskip

\noindent Si richiede di:

\begin{enumerate}
    \item decidere se effettuare l'investimento utilizzando il \textbf{VAN};
    \item determinare come cambierebbe la redditività del progetto se il tasso di sconto aumentasse al \textbf{12\% annuo}.
\end{enumerate}

\subsubsection*{Svolgimento}

\noindent\textbf{1. Calcolo del VAN con $i=10\%$}

\noindent Attualizziamo tutti i flussi di cassa previsti nei cinque anni:

\[
VA=
\frac{30.000}{1,10^1}+
\frac{40.000}{1,10^2}+
\frac{50.000}{1,10^3}+
\frac{60.000}{1,10^4}+
\frac{70.000}{1,10^5}
=182.341,6
\]

\noindent Calcoliamo il VAN considerando l'investimento iniziale di 100.000 euro:

\[
VAN=-100.000+182.341,6=82.341,6
\]

\[
\boxed{VAN\simeq82.342\text{ euro}}
\]

\noindent Poiché il $VAN>0$, il progetto \textbf{crea valore} e l'investimento risulta conveniente.

\medskip
\noindent\textbf{2. Calcolo del VAN con $i=12\%$}

\noindent Aumentiamo il tasso di attualizzazione al 12\% e ricalcoliamo il valore attuale:

\[
VA=
\frac{30.000}{1,12^1}+
\frac{40.000}{1,12^2}+
\frac{50.000}{1,12^3}+
\frac{60.000}{1,12^4}+
\frac{70.000}{1,12^5}
=172.113,4
\]

\noindent Il nuovo VAN è quindi:

\[
VAN=-100.000+172.113,4=72.113,4
\]

\[
\boxed{VAN\simeq72.113\text{ euro}}
\]

\noindent Il VAN rimane positivo, quindi il progetto è ancora \textbf{conveniente}. Tuttavia, aumentando il tasso di attualizzazione dal 10\% al 12\%, il valore attuale dei flussi futuri diminuisce e, 
di conseguenza, diminuisce anche il VAN:

\[
i\uparrow\quad\Rightarrow\quad VA\downarrow\quad\Rightarrow\quad VAN\downarrow
\]
\newpage
\subsection{Esercizio 2}

La società \textbf{Alma Chimica S.p.A.} sta valutando l'acquisto di un nuovo macchinario del costo di \textbf{6,5 milioni di euro}.

\noindent I ricavi previsti ammontano a \textbf{3 milioni di euro all'anno}, mentre i costi di esercizio sono pari a \textbf{2 milioni di euro all'anno}.

\noindent Inoltre:
\begin{itemize}
    \item dopo il quinto anno è necessaria una manutenzione straordinaria del costo di \textbf{1 milione di euro};
    \item dopo il decimo anno è necessaria un'ulteriore manutenzione straordinaria del costo di \textbf{1 milione di euro};
    \item al termine del quindicesimo anno il macchinario può essere venduto per \textbf{800.000 euro}.
\end{itemize}

\noindent Considerando un tasso di interesse del \textbf{7\%}, determinare il \textbf{VAN del macchinario}.

\subsubsection*{Svolgimento}

\noindent Il costo iniziale del macchinario è:

\[
F_0=-6.500.000
\]

\noindent Ogni anno il macchinario genera ricavi per 3 milioni di euro e costi di esercizio per 2 milioni di euro. Il \textbf{flusso di cassa netto annuo} è quindi:

\[
F=3.000.000-2.000.000=1.000.000
\]

\noindent Poiché il flusso di 1 milione di euro si ripete ogni anno per 15 anni, possiamo utilizzare la formula del valore attuale di una \textbf{rendita annua costante}:

\[
VA=C\left[\frac{1}{r}-\frac{1}{r(1+r)^t}\right]
\]

\noindent con:
\begin{itemize}
    \item $C=1.000.000$ $\rightarrow$ flusso di cassa netto annuo;
    \item $r=0,07$ $\rightarrow$ tasso di attualizzazione;
    \item $t=15$ $\rightarrow$ durata dell'investimento.
\end{itemize}

\noindent Sostituendo i valori:

\[
VA=
1.000.000
\left[
\frac{1}{0,07}
-
\frac{1}{0,07(1,07)^{15}}
\right]
=9.107.914
\]

\noindent Dobbiamo poi considerare le due \textbf{manutenzioni straordinarie} da 1 milione di euro, sostenute rispettivamente alla fine del quinto e del decimo anno.

\noindent Essendo delle uscite di denaro, vengono considerate con segno negativo e attualizzate separatamente:

\[
VA_{\text{manutenzioni}}
=
-\frac{1.000.000}{(1,07)^5}
-\frac{1.000.000}{(1,07)^{10}}
\]

\[
VA_{\text{manutenzioni}}
=-1.221.335
\]

\noindent Quindi, fino a questo punto, abbiamo:
\begin{itemize}
    \item $+9.107.914$ euro $\rightarrow$ valore attuale dei flussi di cassa annuali;
    \item $-1.221.335$ euro $\rightarrow$ valore attuale delle manutenzioni straordinarie.
\end{itemize}
\noindent Al termine dei 15 anni il macchinario viene venduto per 800.000 euro.
Poiché questo incasso avverrà al quindicesimo anno, dobbiamo attualizzarlo per
conoscere quanto vale oggi:

\[
VA_{\text{vendita}}
=
\frac{800.000}{(1,07)^{15}}
=
289.956,8
\]

\noindent A questo punto possiamo calcolare il \textbf{VAN dell'intero progetto},
sommando tutti i flussi di cassa attualizzati e considerando il costo iniziale
del macchinario:

\[
VAN
=
-6.500.000
+9.107.914
-1.221.335
+289.956,8
\]

\[
\boxed{VAN=1.676.535,80\text{ euro}}
\]

\noindent Il VAN è \textbf{positivo}, quindi il progetto è conveniente:
considerando anche il valore temporale del denaro, l'investimento genera un
valore aggiuntivo di circa 1,68 milioni di euro.
\newpage
\subsection{Esercizio 3}

\noindent La società Digital Technologies S.r.l. sta valutando due progetti di
investimento alternativi, il \textbf{Progetto X} e il \textbf{Progetto Y},
entrambi con durata di 4 anni e caratterizzati dai seguenti flussi di cassa.

\medskip
\noindent\textbf{Progetto X:}
\begin{itemize}
    \item Investimento iniziale: 600.000 euro
    \item Anno 1: 150.000 euro
    \item Anno 2: 180.000 euro
    \item Anno 3: 210.000 euro
    \item Anno 4: 240.000 euro
\end{itemize}

\noindent\textbf{Progetto Y:}
\begin{itemize}
    \item Investimento iniziale: 750.000 euro
    \item Anno 1: 200.000 euro
    \item Anno 2: 220.000 euro
    \item Anno 3: 240.000 euro
    \item Anno 4: 260.000 euro
\end{itemize}

\noindent Considerando un tasso di sconto del 10\% annuo, si determini quale
progetto di investimento finanziare utilizzando il VAN.
\subsubsection*{Svolgimento}

\noindent\textbf{Progetto X}

\noindent Attualizziamo i quattro flussi di cassa utilizzando il tasso di sconto
del 10\%:

\[
VA_X=
\frac{150.000}{1,10^1}
+\frac{180.000}{1,10^2}
+\frac{210.000}{1,10^3}
+\frac{240.000}{1,10^4}
=606.823,30
\]

\noindent Sottraiamo dal valore attuale dei flussi l'investimento iniziale:

\[
VAN_X=-600.000+606.823,30
=\boxed{6.823,30}
\]

\noindent\textbf{Progetto Y}

\noindent Allo stesso modo attualizziamo i flussi di cassa del secondo progetto:

\[
VA_Y=
\frac{200.000}{1,10^1}
+\frac{220.000}{1,10^2}
+\frac{240.000}{1,10^3}
+\frac{260.000}{1,10^4}
=721.535,41
\]

\noindent Sottraiamo l'investimento iniziale:

\[
VAN_Y=-750.000+721.535,41
=\boxed{-28.464,59}
\]

\noindent Il \textbf{Progetto X è da preferire}, perché presenta un VAN positivo
pari a 6.823,30 euro e quindi genera valore per l'impresa.

\noindent Il Progetto Y presenta invece un VAN negativo: i flussi di cassa
attualizzati non sono sufficienti a recuperare l'investimento iniziale.
\newpage
\subsection{Esercizio 4}

\noindent La società InfoTech S.r.l. sta prendendo in considerazione due progetti
di investimento alternativi, il \textbf{Progetto A} e il \textbf{Progetto B},
entrambi di durata quinquennale e caratterizzati dai seguenti flussi di cassa:

\noindent\textbf{Progetto A:}
\begin{itemize}
    \item Investimento iniziale: 110.000 euro
    \item Flusso di cassa netto annuo: 50.000 euro
\end{itemize}

\noindent\textbf{Progetto B:}
\begin{itemize}
    \item Investimento iniziale: 150.000 euro
    \item Flusso di cassa netto annuo: 60.000 euro
\end{itemize}

\noindent Considerando un tasso di sconto del 10\% annuo, si determini quale
progetto di investimento finanziare utilizzando il VAN.
\subsubsection*{Svolgimento}

\noindent Poiché entrambi i progetti generano un flusso di cassa costante per
5 anni, utilizziamo la formula del valore attuale di una rendita annua costante.

\noindent\textbf{Progetto A}

\[
VA_A=
50.000
\left[
\frac{1}{0,10}
-\frac{1}{0,10(1,10)^5}
\right]
=189.539
\]

\noindent Calcoliamo quindi il VAN sottraendo l'investimento iniziale:

\[
VAN_A=-110.000+189.539
=\boxed{79.539}
\]

\noindent\textbf{Progetto B}

\[
VA_B=
60.000
\left[
\frac{1}{0,10}
-\frac{1}{0,10(1,10)^5}
\right]
=227.447
\]

\noindent Calcoliamo il VAN:

\[
VAN_B=-150.000+227.447
=\boxed{77.447}
\]

\noindent Entrambi i progetti hanno un VAN positivo e sono quindi convenienti.
Tuttavia, essendo progetti alternativi, viene scelto il \textbf{Progetto A},
poiché presenta il VAN maggiore:

\[
\boxed{VAN_A=79.539 > VAN_B=77.447}
\]

\noindent Anche se il Progetto B genera flussi di cassa annuali maggiori,
richiede anche un investimento iniziale più elevato. Il VAN tiene conto di
entrambi gli aspetti e mostra che il Progetto A crea leggermente più valore.
\newpage
\section{Esercitazione -- Analisi degli investimenti}

\subsection*{Riepilogo delle formule}

\noindent Prima di svolgere gli esercizi, richiamiamo le principali formule
utilizzate per valutare la convenienza economica degli investimenti.

\subsubsection*{Valore Attuale Netto (VAN)}

\noindent Il VAN rappresenta il valore creato da un investimento,
attualizzando tutti i flussi di cassa al costo opportunità del capitale.

\[
\boxed{
VAN=F_0+\sum_{t=1}^{N}\frac{F_t}{(1+i)^t}
}
\]

\noindent Un progetto è conveniente quando $VAN>0$.

\subsubsection*{Tasso Interno di Rendimento (TIR)}

\noindent Il TIR è il tasso di attualizzazione che rende il VAN uguale a zero.

\[
\boxed{
0=F_0+\sum_{t=1}^{N}\frac{F_t}{(1+TIR)^t}
}
\]

\noindent Per investimenti tradizionali, un progetto è conveniente
quando $TIR>i$, dove $i$ rappresenta il costo opportunità del capitale.

\subsubsection*{Indice di redditività}

\noindent L'indice di redditività permette di confrontare il valore
creato rispetto al capitale investito.

\[
\boxed{
PI=\frac{VAN}{I_0}
}
\]

\noindent Dove $I_0$ rappresenta l'investimento iniziale.
Un valore maggiore indica una maggiore creazione di valore
per unità di capitale investito.

\subsubsection*{Razionamento del capitale}

\noindent Quando il budget è limitato, bisogna individuare
la combinazione di investimenti che massimizza il VAN complessivo.

\[
\boxed{
\max \sum_{j=1}^{m} VAN_j
}
\]

\noindent rispettando il vincolo:

\[
\boxed{
\sum_{j=1}^{m} I_{0,j}\leq B
}
\]

\noindent dove $B$ è il budget disponibile.

\subsubsection*{Payback Period}

\noindent Il Payback Period indica il tempo necessario
per recuperare l'investimento iniziale.

\noindent Se il recupero avviene durante un anno:

\[
\boxed{
PBP=k+\frac{C_{\text{mancante}}}{F_{k+1}}
}
\]

\noindent dove $k$ indica gli anni già trascorsi,
$C_{\text{mancante}}$ il capitale ancora da recuperare
e $F_{k+1}$ il flusso dell'anno successivo.

\noindent Nel caso di flussi annui costanti:

\[
\boxed{
PBP=\frac{I_0}{R}
}
\]

\subsubsection*{Payback Period attualizzato}

\noindent Il Payback attualizzato considera il valore temporale
del denaro, utilizzando i flussi di cassa attualizzati:

\[
\boxed{
F_t^{\text{att}}=\frac{F_t}{(1+i)^t}
}
\]

\noindent Il recupero avviene quando la somma dei flussi
attualizzati raggiunge l'investimento iniziale.

\subsubsection*{Valore attuale di una rendita costante}

\noindent Quando un investimento genera entrate annue costanti,
si utilizza il fattore di attualizzazione della rendita:

\[
\boxed{
FA(i,n)=\frac{1-(1+i)^{-n}}{i}
}
\]

\noindent Il VAN diventa quindi:

\[
\boxed{
VAN=-I_0+R\cdot FA(i,n)
}
\]

\noindent dove $R$ rappresenta il flusso annuo costante.


\subsection{Esercizio 1 -- Valutazione tramite VAN e TIR}

\noindent Un progetto presenta i seguenti flussi di cassa:

\[
F_0=+9\,500
\qquad
F_1=+9\,000
\qquad
F_2=-21\,000
\]

\noindent Il costo opportunità del capitale è pari all'8\%.
Si vuole valutare la convenienza dell'investimento.

\subsubsection*{Calcolo del TIR}

\[
0=9\,500+\frac{9\,000}{1+TIR}
-\frac{21\,000}{(1+TIR)^2}
\]

\noindent Ponendo $x=1+TIR$:

\[
9\,500x^2+9\,000x-21\,000=0
\]

\noindent Dividendo per 100:

\[
95x^2+90x-210=0
\]

\noindent Risolvendo:

\[
x_1\simeq1{,}0867
\qquad
x_2\simeq-2{,}0341
\]

\noindent Consideriamo solamente la soluzione ammissibile $x>0$:

\[
TIR=1{,}0867-1
\]

\[
\boxed{TIR\simeq8{,}67\%}
\]

\subsubsection*{Calcolo del VAN}

\[
VAN=9\,500+\frac{9\,000}{1{,}08}
-\frac{21\,000}{1{,}08^2}
\]

\[
\boxed{VAN\simeq-170{,}78}
\]

\noindent \textbf{Conclusione:} il progetto non è conveniente,
poiché presenta un VAN negativo.

\medskip

\noindent Nonostante il TIR sia maggiore dell'8\%, i flussi
non rappresentano un investimento tradizionale,
poiché si ricevono entrate iniziali e si sostiene un'uscita finale.

\noindent In questi casi il criterio del TIR può essere fuorviante,
quindi si preferisce utilizzare il VAN.


\subsection{Esercizio 2 -- Razionamento del capitale}

\noindent Si dispone di un budget massimo di $90\,000$
da investire nei seguenti progetti:

\begin{center}
\begin{tabular}{|c|r|r|c|}
\hline
\textbf{Progetto} & \textbf{VAN} &
\textbf{Investimento} & \textbf{PI} \\
\hline
1 & 5.000 & 10.000 & 0,50 \\
2 & 5.000 & 5.000 & 1,00 \\
3 & 10.000 & 90.000 & 0,11 \\
4 & 15.000 & 60.000 & 0,25 \\
5 & 15.000 & 75.000 & 0,20 \\
6 & 3.000 & 15.000 & 0,20 \\
\hline
\end{tabular}
\end{center}

\noindent L'obiettivo è individuare la combinazione di progetti
che massimizza il VAN complessivo, senza superare il budget.

\subsubsection*{Calcolo dell'indice di redditività}

\noindent Per ogni progetto calcoliamo:

\[
PI=\frac{VAN}{I_0}
\]

\noindent Ad esempio, per il progetto 4:

\[
PI_4=\frac{15\,000}{60\,000}=0{,}25
\]

\noindent Considerando gli indici di redditività e
verificando le combinazioni ammissibili,
la scelta ottimale è costituita dai progetti
\textbf{1, 2, 4 e 6}.

\subsubsection*{Verifica del budget}

\[
I_{\text{totale}}
=10\,000+5\,000+60\,000+15\,000
\]

\[
\boxed{I_{\text{totale}}=90\,000}
\]

\noindent Il vincolo di budget viene rispettato.

\subsubsection*{Calcolo del VAN complessivo}

\[
VAN_{\text{totale}}
=5\,000+5\,000+15\,000+3\,000
\]

\[
\boxed{VAN_{\text{totale}}=28\,000}
\]

\noindent \textbf{Conclusione:} scegliamo i progetti
\textbf{1, 2, 4 e 6}, poiché permettono di ottenere
il massimo VAN complessivo rispettando il capitale disponibile.

\medskip

\noindent Quando i progetti sono indivisibili,
il semplice ordinamento tramite PI non garantisce
sempre la combinazione ottimale.


\subsection{Esercizio 3 -- VAN e tempo di recupero}

\noindent Si considerano tre progetti con un costo opportunità
del capitale pari al $6\%$.

\begin{center}
\small
\begin{tabular}{|c|r|r|r|r|r|r|}
\hline
\textbf{Progetto} & $F_0$ & $F_1$ & $F_2$ &
$F_3$ & $F_4$ & $F_5$ \\
\hline
1 & -1000 & 1200 & 2500 & 2000 & 1750 & 1500 \\
2 & -2000 & 600 & 1000 & 2000 & 3000 & 5000 \\
3 & -3000 & 4000 & 200 & 0 & 1000 & 1500 \\
\hline
\end{tabular}
\end{center}

\subsubsection*{Calcolo del VAN}

\noindent Per il progetto 1:

\[
\begin{aligned}
VAN_1={}&-1000+\frac{1200}{1{,}06}
+\frac{2500}{1{,}06^2}\\
&+\frac{2000}{1{,}06^3}
+\frac{1750}{1{,}06^4}
+\frac{1500}{1{,}06^5}
\end{aligned}
\]

\[
\boxed{VAN_1\simeq6\,543{,}36}
\]

\noindent Per il progetto 2:

\[
\begin{aligned}
VAN_2={}&-2000+\frac{600}{1{,}06}
+\frac{1000}{1{,}06^2}\\
&+\frac{2000}{1{,}06^3}
+\frac{3000}{1{,}06^4}
+\frac{5000}{1{,}06^5}
\end{aligned}
\]

\[
\boxed{VAN_2\simeq7\,247{,}84}
\]

\noindent Per il progetto 3:

\[
\begin{aligned}
VAN_3={}&-3000+\frac{4000}{1{,}06}
+\frac{200}{1{,}06^2}\\
&+\frac{1000}{1{,}06^4}
+\frac{1500}{1{,}06^5}
\end{aligned}
\]

\[
\boxed{VAN_3\simeq2\,864{,}57}
\]

\subsubsection*{Calcolo del Payback Period}

\noindent \textbf{Progetto 1}

\noindent L'investimento viene recuperato durante il primo anno:

\[
PBP_1=\frac{1000}{1200}
\]

\[
\boxed{PBP_1\simeq0{,}83\text{ anni}}
\]

\medskip

\noindent \textbf{Progetto 2}

\noindent Dopo due anni sono stati recuperati:

\[
600+1000=1600
\]

\noindent Il capitale ancora da recuperare è:

\[
2000-1600=400
\]

\[
PBP_2=2+\frac{400}{2000}
\]

\[
\boxed{PBP_2=2{,}20\text{ anni}}
\]

\medskip

\noindent \textbf{Progetto 3}

\noindent Il capitale iniziale viene recuperato
durante il primo anno:

\[
PBP_3=\frac{3000}{4000}
\]

\[
\boxed{PBP_3=0{,}75\text{ anni}}
\]

\subsubsection*{Conclusione}

\begin{center}
\begin{tabular}{|c|r|c|}
\hline
\textbf{Progetto} & \textbf{VAN} & \textbf{PBP} \\
\hline
1 & 6.543,36 & 0,83 anni \\
2 & 7.247,84 & 2,20 anni \\
3 & 2.864,57 & 0,75 anni \\
\hline
\end{tabular}
\end{center}

\noindent Se il \textbf{cutoff period è pari a 2 anni},
vengono accettati i progetti il cui tempo di recupero
non supera tale limite.

\[
\boxed{\text{Progetti accettati: 1 e 3}}
\]

\noindent Il progetto 2 viene escluso dal criterio del Payback,
nonostante presenti il VAN maggiore.

\medskip

\noindent \textbf{Nota:} considerando solamente gli anni interi,
i tempi di recupero diventano rispettivamente 1, 3 e 1 anno.
La decisione finale rimane invariata.


\subsection{Esercizio 4 -- Indice di redditività e Payback attualizzato}

\noindent Si considerano quattro progetti alternativi
della durata di 6 anni, con un costo opportunità
del capitale pari al $12\%$.

\begin{center}
\begin{tabular}{|c|r|r|}
\hline
\textbf{Progetto} & \textbf{Investimento} &
\textbf{Entrata annua} \\
\hline
1 & 14.000 & 3.500 \\
2 & 11.000 & 2.800 \\
3 & 13.000 & 3.100 \\
4 & 17.000 & 4.400 \\
\hline
\end{tabular}
\end{center}

\subsubsection*{a) Scelta tramite indice di redditività}

\noindent Poiché le entrate sono costanti, utilizziamo
il fattore di attualizzazione della rendita:

\[
FA(12\%,6)
=\frac{1-(1{,}12)^{-6}}{0{,}12}
\]

\[
\boxed{FA\simeq4{,}1114}
\]

\noindent Calcoliamo il VAN di ciascun progetto:

\[
VAN_1=-14\,000+3\,500\cdot4{,}1114
\]

\[
\boxed{VAN_1\simeq389{,}93}
\]

\[
VAN_2=-11\,000+2\,800\cdot4{,}1114
\]

\[
\boxed{VAN_2\simeq511{,}94}
\]

\[
VAN_3=-13\,000+3\,100\cdot4{,}1114
\]

\[
\boxed{VAN_3\simeq-254{,}64}
\]

\[
VAN_4=-17\,000+4\,400\cdot4{,}1114
\]

\[
\boxed{VAN_4\simeq1\,090{,}19}
\]

\noindent Calcoliamo ora gli indici di redditività:

\[
PI_1=\frac{389{,}93}{14\,000}\simeq0{,}0279
\]

\[
PI_2=\frac{511{,}94}{11\,000}\simeq0{,}0465
\]

\[
PI_3=\frac{-254{,}64}{13\,000}\simeq-0{,}0196
\]

\[
PI_4=\frac{1\,090{,}19}{17\,000}\simeq0{,}0641
\]

\noindent Il progetto 4 presenta l'indice maggiore:

\[
\boxed{PI_4\simeq0{,}0641}
\]

\noindent \textbf{Conclusione:} scegliamo il
\textbf{progetto 4}, che presenta anche il VAN maggiore.

\subsubsection*{b) Tempo di recupero del progetto 4}

\noindent Consideriamo:

\[
I_0=17\,000
\qquad
R=4\,400
\qquad
i=12\%
\]

\noindent \textbf{Payback non attualizzato}

\[
PBP=\frac{17\,000}{4\,400}
\]

\[
PBP\simeq3{,}8636\text{ anni}
\]

\noindent Convertiamo la parte decimale in mesi e giorni,
considerando mesi di 30 giorni:

\[
0{,}8636\cdot12\simeq10{,}36\text{ mesi}
\]

\[
0{,}36\cdot30\simeq11\text{ giorni}
\]

\[
\boxed{
PBP\simeq3\text{ anni, }10\text{ mesi e }11\text{ giorni}
}
\]

\medskip

\noindent \textbf{Payback attualizzato}

\noindent Attualizziamo le entrate annuali:

\[
F_t^{\text{att}}=\frac{4\,400}{(1{,}12)^t}
\]

\noindent Dopo cinque anni il capitale non è ancora
completamente recuperato, mentre viene recuperato
durante il sesto anno.

\medskip

\noindent Calcoliamo la frazione dell'ultimo anno necessaria
per recuperare il capitale residuo:

\[
f=
\frac{
17\,000-\displaystyle\sum_{t=1}^{5}
\frac{4\,400}{(1{,}12)^t}
}{
\frac{4\,400}{(1{,}12)^6}
}
\]

\[
f\simeq0{,}5109
\]

\noindent Pertanto, interpolando il flusso del sesto anno:

\[
PBP_{\text{att}}\simeq5{,}5109\text{ anni}
\]

\noindent Convertendo in mesi e giorni:

\[
0{,}5109\cdot12\simeq6{,}13\text{ mesi}
\]

\[
0{,}13\cdot30\simeq4\text{ giorni}
\]

\[
\boxed{
PBP_{\text{att}}\simeq
5\text{ anni, }6\text{ mesi e }4\text{ giorni}
}
\]

\subsubsection*{Conclusione}

\noindent Il progetto selezionato è il \textbf{progetto 4}.

\[
\boxed{
\begin{aligned}
VAN &\simeq1\,090{,}19\\
PI &\simeq0{,}0641\\
PBP &\simeq3{,}86\text{ anni}\\
PBP_{\text{att}} &\simeq5{,}51\text{ anni}
\end{aligned}
}
\]

\noindent Il Payback attualizzato è maggiore di quello semplice,
poiché considera il valore temporale del denaro.
\newpage
\section{Esempio Esonero}

\subsection{Breve richiamo alla teoria e alle formule}

\subsubsection*{Attualizzazione}

\noindent Attualizzare significa riportare
una somma futura al suo valore di oggi:

\[
\boxed{
VA=\frac{F_t}{(1+i)^t}
}
\]

\begin{itemize}
    \item $VA$: valore attuale;
    \item $F_t$: flusso disponibile al tempo $t$;
    \item $i$: tasso per periodo;
    \item $t$: numero di periodi
    che separano il flusso da oggi.
\end{itemize}

\noindent Con un tasso positivo,
un'entrata futura vale oggi meno
del suo importo nominale.


\subsubsection*{Valore Attuale Netto}

\noindent Il \textbf{Valore Attuale Netto},
indicato con $VAN$,
misura il valore attuale delle entrate
al netto dell'investimento iniziale:

\[
\boxed{
VAN=-I_0+\sum_{t=1}^{n}
\frac{F_t}{(1+i)^t}
}
\]

\begin{itemize}
    \item $I_0$: investimento iniziale,
    espresso come importo positivo;
    \item $F_t$: flusso netto del periodo $t$;
    \item $n$: numero complessivo dei periodi.
\end{itemize}

\noindent Il segno meno davanti a $I_0$
indica un'uscita iniziale di denaro.

\[
\begin{aligned}
VAN>0
&\Rightarrow \text{il progetto crea valore};\\
VAN=0
&\Rightarrow \text{remunera esattamente il tasso richiesto};\\
VAN<0
&\Rightarrow \text{non remunera il tasso richiesto}.
\end{aligned}
\]


\subsubsection*{Tasso Interno di Rendimento}

\noindent Il \textbf{TIR}
è il tasso che rende nullo il $VAN$:

\[
\boxed{
0=-I_0+\sum_{t=1}^{n}
\frac{F_t}{(1+TIR)^t}
}
\]

\noindent Per questi progetti,
con un'uscita iniziale
seguita da entrate positive:

\[
\boxed{
TIR>i
\Rightarrow
\text{progetto accettabile}
}
\]

\noindent \textbf{In questo esercizio
non serve trovare il TIR preciso.}

\noindent Ci basta sapere
se supera il tasso richiesto del 10\%.

\[
\boxed{
VAN(10\%)>0
\Longleftrightarrow
TIR>10\%
}
\]

\noindent Infatti,
aumentando il tasso diminuisce
il valore attuale delle entrate.

\noindent Se al 10\% il $VAN$
è ancora positivo,
il tasso che lo porta a zero
deve essere più alto.


\subsubsection*{Rendite costanti}

\noindent Una \textbf{rendita costante}
è una successione di pagamenti
dello stesso importo,
effettuati a intervalli regolari.

\noindent Se i pagamenti vengono ricevuti
alla fine di ogni periodo,
la rendita è \textbf{posticipata}.

\begin{itemize}
    \item $R$: importo di ciascun pagamento;
    \item $i$: tasso per periodo;
    \item $n$: numero dei pagamenti.
\end{itemize}

\noindent Il tasso deve riferirsi
allo stesso intervallo dei pagamenti:
con pagamenti annuali
utilizziamo un tasso annuo.


\subsubsection*{Coefficiente di attualizzazione:
il CAF}

\noindent Il valore attuale
dei pagamenti costanti è:

\[
VA=
\frac{R}{1+i}
+
\frac{R}{(1+i)^2}
+
\cdots
+
\frac{R}{(1+i)^n}.
\]

\noindent Raccogliendo il pagamento $R$:

\[
VA=
R\left[
\frac{1}{1+i}
+
\frac{1}{(1+i)^2}
+
\cdots
+
\frac{1}{(1+i)^n}
\right].
\]

\noindent La quantità tra parentesi
è il \textbf{coefficiente di attualizzazione
di una rendita costante},
indicato con $CAF$:

\[
\boxed{
CAF(i,n)
=
\sum_{t=1}^{n}\frac{1}{(1+i)^t}
}
\]

\noindent La somma si può scrivere:

\[
CAF(i,n)
=
\frac{1-(1+i)^{-n}}{i}.
\]

\noindent Separando i termini:

\[
\boxed{
CAF(i,n)
=
\frac{1}{i}
-
\frac{1}{i(1+i)^n}
}
\]

\noindent Quindi:

\[
\boxed{
VA=CAF(i,n)\cdot R
}
\]

\noindent Se il pagamento è indicato con $C$:

\[
\boxed{
VA=C\cdot CAF(i,n)
}
\]

\noindent Cambia soltanto la lettera
utilizzata per il pagamento.


\subsubsection*{Collegamento con la notazione degli appunti}

\noindent Negli appunti il tasso
può essere indicato con $r$
e il numero di periodi con $t$:

\[
\boxed{
i=r,
\qquad
n=t
}
\]

\noindent Pertanto:

\[
CAF(r,t)
=
\frac{1}{r}
-
\frac{1}{r(1+r)^t}.
\]

\noindent In questa scrittura $t$
indica il numero totale dei pagamenti.


\subsubsection*{VAN con rendite costanti}

\noindent Sottraiamo l'investimento iniziale
dal valore attuale delle entrate:

\[
\boxed{
VAN=-I_0+CAF(i,n)\cdot R
}
\]

\begin{enumerate}
    \item Calcoliamo il $CAF$.
    \item Lo moltiplichiamo per il flusso annuo.
    \item Sottraiamo l'investimento iniziale.
\end{enumerate}


\subsubsection*{Esempio: tasso del 10\% e cinque pagamenti}

\[
i=10\%=0{,}10,
\qquad
n=5.
\]

\[
CAF(10\%,5)
=
\frac{1}{0{,}10}
-
\frac{1}{0{,}10(1{,}10)^5}.
\]

\noindent Il primo termine è:

\[
\frac{1}{0{,}10}=10{,}00.
\]

\noindent Il secondo termine vale circa:

\[
\frac{1}{0{,}10(1{,}10)^5}
\simeq6{,}21.
\]

\noindent Quindi:

\[
CAF(10\%,5)
\simeq10{,}00-6{,}21
=
\boxed{3{,}79}.
\]

\noindent Cinque pagamenti annuali di 1
valgono oggi circa 3,79,
con un tasso del 10\%.

\medskip

\noindent \textbf{Convenzione utilizzata:}
nei punti a, b e d utilizziamo
direttamente il coefficiente arrotondato
dello svolgimento del professore:

\[
\boxed{CAF=3{,}79}
\]

\noindent I VAN ottenuti sono quindi
approssimati con questo coefficiente.


\subsubsection*{VAN aggiornato con una spesa aggiuntiva}

\noindent Una spesa aggiuntiva
riduce il VAN del progetto.

\noindent Prima attualizziamo la spesa,
poi la sottraiamo al VAN precedente.

\noindent Se la spesa si verifica
alla fine del primo anno:

\[
\boxed{
VAN_{\text{aggiornato}}
=
VAN_{\text{precedente}}
-
\frac{\text{Spesa}}{(1+i)^1}
}
\]

\begin{enumerate}
    \item Riportiamo la spesa futura a oggi.
    \item La sottraiamo al VAN precedente.
    \item Controlliamo se il VAN aggiornato
    è ancora positivo.
\end{enumerate}

\[
\boxed{
VAN_{\text{aggiornato}}>0
\Rightarrow
\text{progetto ancora accettabile}
}
\]

\noindent Il confronto è con zero,
non con la soglia del punto b.


\subsubsection*{Indice di redditività:
il PI}

\noindent Seguendo la definizione
utilizzata dal professore,
l'indice di redditività è:

\[
\boxed{
PI=\frac{VAN}{I_0}
}
\]

\noindent Il $PI$ misura il valore attuale
netto creato per ogni unità
di investimento iniziale.

\noindent Per esprimerlo in percentuale
moltiplichiamo il rapporto per 100.

\medskip

\noindent Per esempio,
un $PI$ del 20\% significa che,
per ogni 100 investiti inizialmente,
il progetto genera un VAN di 20.

\noindent \textbf{Attenzione:}
il PI così definito
non è un rendimento annuo
e non coincide con il TIR.


\subsubsection*{PI di un paniere:
media ponderata}

\noindent Per un paniere di progetti indipendenti,
sommiamo i VAN e gli investimenti:

\[
\boxed{
PI_{\text{paniere}}
=
\frac{\sum_j VAN_j}{\sum_j I_{0,j}}
}
\]

\noindent Poiché:

\[
VAN_j=PI_j\cdot I_{0,j},
\]

\noindent possiamo scrivere:

\[
\boxed{
PI_{\text{paniere}}
=
\frac{
\sum_j PI_j\cdot I_{0,j}
}{
\sum_j I_{0,j}
}
}
\]

\noindent È una \textbf{media ponderata}
dei PI dei singoli progetti,
con gli investimenti iniziali come pesi.

\noindent Un progetto che richiede
un investimento maggiore
pesa di più nella media.

\medskip

\noindent Per un paniere formato da B e F:

\[
\boxed{
PI_{B+F}
=
\frac{
PI_B\cdot I_{0,B}
+
PI_F\cdot I_{0,F}
}{
I_{0,B}+I_{0,F}
}
}
\]

\noindent Non utilizziamo la media semplice,
perché B e F richiedono investimenti diversi.

\noindent Nei calcoli dei pesi utilizziamo
i PI come rapporti;
i risultati vengono poi espressi in percentuale.


\subsubsection*{Recupero attualizzato:
saldo cumulato}

\noindent Per verificare quando recuperiamo
l'investimento,
partiamo dall'uscita iniziale
e aggiungiamo ogni anno
il flusso attualizzato.

\noindent Chiamiamo $S_t$
il \textbf{saldo cumulato attualizzato}
alla fine dell'anno $t$.

\noindent Al tempo iniziale:

\[
\boxed{S_0=-I_0}
\]

\noindent Alla fine del primo anno:

\[
S_1=-I_0+\frac{R}{1+i}.
\]

\noindent Alla fine del secondo anno:

\[
S_2=S_1+\frac{R}{(1+i)^2}.
\]

\noindent Alla fine del terzo anno:

\[
S_3=S_2+\frac{R}{(1+i)^3}.
\]

\noindent In generale:

\[
\boxed{
S_t=S_{t-1}+\frac{R}{(1+i)^t}
}
\]

\noindent Equivalentemente:

\[
\boxed{
S_t=-I_0+\sum_{k=1}^{t}
\frac{R}{(1+i)^k}
}
\]

\noindent Il significato è semplice:

\begin{itemize}
    \item se $S_t<0$,
    dobbiamo ancora recuperare
    una parte dell'investimento;

    \item se $S_t=0$,
    abbiamo recuperato esattamente
    l'investimento attualizzato;

    \item se $S_t>0$,
    l'investimento è stato recuperato
    e rimane un margine positivo.
\end{itemize}

\noindent Per recuperare entro tre anni:

\[
\boxed{S_3\geq0}
\]

\noindent Non basta sommare
i flussi nominali:
ogni entrata deve prima
essere attualizzata.


\subsection{Testo e dati dell'esercizio}

\noindent Una società informatica valuta
sei progetti di cybersecurity.

\noindent Ogni progetto presenta:

\begin{itemize}
    \item durata di cinque anni;
    \item valore di recupero finale nullo;
    \item investimento iniziale al tempo zero;
    \item flussi annui costanti.
\end{itemize}

\noindent Interpretiamo la redditività annua
come il flusso netto ricevuto
alla fine di ogni anno.

\begin{center}
    \begin{tabular}{|c|r|r|}
        \hline
        \textbf{Progetto}
        & \textbf{Investimento iniziale}
        & \textbf{Flusso annuo} \\
        \hline
        A & 80.000,00 & 12.500,00 \\
        B & 90.000,00 & 27.500,00 \\
        C & 50.000,00 & 5.000,00 \\
        D & 65.000,00 & 10.500,00 \\
        E & 100.000,00 & 35.000,00 \\
        F & 55.000,00 & 24.000,00 \\
        \hline
    \end{tabular}
\end{center}

\noindent Il tasso richiesto è:

\[
\boxed{i=10\%=0{,}10}
\]

\noindent Non aggiungiamo un ricavo finale,
perché il valore di recupero è nullo.


\subsection{Punto a:
progetti accettabili secondo il TIR}

\subsubsection*{Che cosa viene richiesto}

\noindent I progetti sono indipendenti
e le risorse finanziarie sono illimitate.

\noindent Possiamo accettare
tutti i progetti con:

\[
\boxed{TIR>10\%}
\]

\noindent Non dobbiamo scegliere
necessariamente un solo progetto.


\subsubsection*{1. Progetto A:
formula completa del TIR}

\noindent Il progetto A richiede 80.000,00
e genera cinque entrate di 12.500,00.

\noindent L'equazione del TIR è:

\[
\begin{aligned}
0={}&-80.000{,}00
+\frac{12.500{,}00}{1+TIR_A}
+\frac{12.500{,}00}{(1+TIR_A)^2}\\
&+\frac{12.500{,}00}{(1+TIR_A)^3}
+\frac{12.500{,}00}{(1+TIR_A)^4}
+\frac{12.500{,}00}{(1+TIR_A)^5}.
\end{aligned}
\]

\noindent Non cerchiamo il TIR preciso.

\noindent Sostituiamo il tasso del 10\%
per controllare il segno del VAN:

\[
\begin{aligned}
VAN_A={}&-80.000{,}00
+\frac{12.500{,}00}{1{,}10}
+\frac{12.500{,}00}{(1{,}10)^2}\\
&+\frac{12.500{,}00}{(1{,}10)^3}
+\frac{12.500{,}00}{(1{,}10)^4}
+\frac{12.500{,}00}{(1{,}10)^5}.
\end{aligned}
\]

\noindent Raccogliendo 12.500,00,
riconosciamo il CAF:

\[
VAN_A
=
-80.000{,}00
+
12.500{,}00\cdot CAF(10\%,5).
\]

\noindent La formula completa
e quella con il CAF sono equivalenti.

\noindent Utilizzando il CAF arrotondato:

\[
VAN_A
\simeq
-80.000{,}00
+
3{,}79\cdot12.500{,}00.
\]

\[
3{,}79\cdot12.500{,}00
=
47.375{,}00.
\]

\[
VAN_A
\simeq
-80.000{,}00+47.375{,}00
=
\boxed{-32.625{,}00}.
\]

\[
\boxed{
TIR_A<10\%
\Rightarrow
\text{A non accettabile}
}
\]


\subsubsection*{2. Progetto B}

\noindent Utilizziamo direttamente:

\[
VAN=-I_0+CAF(i,n)\cdot R.
\]

\[
VAN_B
\simeq
-90.000{,}00+3{,}79\cdot27.500{,}00.
\]

\[
3{,}79\cdot27.500{,}00
=
104.225{,}00.
\]

\[
VAN_B
\simeq
-90.000{,}00+104.225{,}00
=
\boxed{14.225{,}00}.
\]

\[
\boxed{
TIR_B>10\%
\Rightarrow
\text{B accettabile}
}
\]


\subsubsection*{3. Progetto C}

\[
VAN_C
\simeq
-50.000{,}00+3{,}79\cdot5.000{,}00.
\]

\[
3{,}79\cdot5.000{,}00
=
18.950{,}00.
\]

\[
VAN_C
\simeq
-50.000{,}00+18.950{,}00
=
\boxed{-31.050{,}00}.
\]

\[
\boxed{
TIR_C<10\%
\Rightarrow
\text{C non accettabile}
}
\]


\subsubsection*{4. Progetto D}

\[
VAN_D
\simeq
-65.000{,}00+3{,}79\cdot10.500{,}00.
\]

\[
3{,}79\cdot10.500{,}00
=
39.795{,}00.
\]

\[
VAN_D
\simeq
-65.000{,}00+39.795{,}00
=
\boxed{-25.205{,}00}.
\]

\[
\boxed{
TIR_D<10\%
\Rightarrow
\text{D non accettabile}
}
\]


\subsubsection*{5. Progetto E}

\[
VAN_E
\simeq
-100.000{,}00+3{,}79\cdot35.000{,}00.
\]

\[
3{,}79\cdot35.000{,}00
=
132.650{,}00.
\]

\[
VAN_E
\simeq
-100.000{,}00+132.650{,}00
=
\boxed{32.650{,}00}.
\]

\[
\boxed{
TIR_E>10\%
\Rightarrow
\text{E accettabile}
}
\]


\subsubsection*{6. Progetto F}

\[
VAN_F
\simeq
-55.000{,}00+3{,}79\cdot24.000{,}00.
\]

\[
3{,}79\cdot24.000{,}00
=
90.960{,}00.
\]

\[
VAN_F
\simeq
-55.000{,}00+90.960{,}00
=
\boxed{35.960{,}00}.
\]

\[
\boxed{
TIR_F>10\%
\Rightarrow
\text{F accettabile}
}
\]


\subsubsection*{Risposta al punto a}

\begin{center}
    \begin{tabular}{|c|r|c|}
        \hline
        \textbf{Progetto}
        & \textbf{VAN approssimato}
        & \textbf{Accettabile} \\
        \hline
        A & -32.625,00 & No \\
        B & 14.225,00 & Sì \\
        C & -31.050,00 & No \\
        D & -25.205,00 & No \\
        E & 32.650,00 & Sì \\
        F & 35.960,00 & Sì \\
        \hline
    \end{tabular}
\end{center}

\[
\boxed{
\text{Progetti accettabili: B, E e F}
}
\]

\noindent Abbiamo controllato
se il TIR supera il 10\%,
senza calcolarne il valore preciso.


\subsection{Punto b:
valore generato di almeno 30.000,00}

\subsubsection*{Che cosa viene richiesto}

\noindent Consideriamo soltanto
i progetti accettati al punto a:

\[
B,\qquad E,\qquad F.
\]

\noindent Seguendo lo svolgimento
del professore,
verifichiamo il requisito sul VAN:

\[
\boxed{VAN\geq30.000{,}00}
\]

\noindent Non calcoliamo
il valore futuro netto.

\noindent Utilizziamo:

\[
VAN=-I_0+CAF(i,n)\cdot R,
\]

\noindent con:

\[
CAF(10\%,5)\simeq3{,}79.
\]


\subsubsection*{1. Progetto B}

\[
I_{0,B}=90.000{,}00,
\qquad
R_B=27.500{,}00.
\]

\[
VAN_B
=
-I_{0,B}+CAF(10\%,5)\cdot R_B.
\]

\[
VAN_B
\simeq
-90.000{,}00+3{,}79\cdot27.500{,}00.
\]

\[
3{,}79\cdot27.500{,}00
=
104.225{,}00.
\]

\[
VAN_B
\simeq
-90.000{,}00+104.225{,}00
=
\boxed{14.225{,}00}.
\]

\noindent Confrontiamo con la soglia:

\[
14.225{,}00<30.000{,}00.
\]

\[
\boxed{
\text{B non soddisfa il requisito}
}
\]


\subsubsection*{2. Progetto E}

\[
I_{0,E}=100.000{,}00,
\qquad
R_E=35.000{,}00.
\]

\[
VAN_E
=
-I_{0,E}+CAF(10\%,5)\cdot R_E.
\]

\[
VAN_E
\simeq
-100.000{,}00+3{,}79\cdot35.000{,}00.
\]

\[
3{,}79\cdot35.000{,}00
=
132.650{,}00.
\]

\[
VAN_E
\simeq
-100.000{,}00+132.650{,}00
=
\boxed{32.650{,}00}.
\]

\noindent Confrontiamo con la soglia:

\[
32.650{,}00>30.000{,}00.
\]

\[
\boxed{
\text{E soddisfa il requisito}
}
\]


\subsubsection*{3. Progetto F}

\[
I_{0,F}=55.000{,}00,
\qquad
R_F=24.000{,}00.
\]

\[
VAN_F
=
-I_{0,F}+CAF(10\%,5)\cdot R_F.
\]

\[
VAN_F
\simeq
-55.000{,}00+3{,}79\cdot24.000{,}00.
\]

\[
3{,}79\cdot24.000{,}00
=
90.960{,}00.
\]

\[
VAN_F
\simeq
-55.000{,}00+90.960{,}00
=
\boxed{35.960{,}00}.
\]

\noindent Confrontiamo con la soglia:

\[
35.960{,}00>30.000{,}00.
\]

\[
\boxed{
\text{F soddisfa il requisito}
}
\]


\subsubsection*{Risposta al punto b}

\[
\boxed{
\text{Progetti accettabili: E e F}
}
\]

\noindent B era accettabile al punto a,
ma non raggiunge
la nuova soglia di 30.000,00.


\subsection{Punto c:
spesa aggiuntiva di 10.000,00 al primo anno}

\subsubsection*{Che cosa viene richiesto}

\noindent Consideriamo i progetti
risultati idonei al punto b:

\[
E,\qquad F.
\]

\noindent Verifichiamo se rimangono accettabili
dopo una spesa aggiuntiva di 10.000,00
alla fine del primo anno.

\noindent Utilizziamo i VAN del punto b:

\[
VAN_E\simeq32.650{,}00,
\qquad
VAN_F\simeq35.960{,}00.
\]

\noindent Aggiorniamo il VAN
dell'intero progetto:

\[
\boxed{
VAN_{\text{aggiornato}}
=
VAN_{\text{precedente}}
-
\frac{\text{Spesa}}{(1+i)^1}
}
\]


\subsubsection*{1. Attualizziamo la spesa}

\[
VA_{\text{spesa}}
=
\frac{10.000{,}00}{(1+0{,}10)^1}.
\]

\[
VA_{\text{spesa}}
=
\frac{10.000{,}00}{1{,}10}
\simeq
\boxed{9.090{,}91}.
\]

\noindent Sottraiamo questo importo,
perché sia la spesa attualizzata
sia il VAN sono riferiti a oggi.


\subsubsection*{2. Progetto E}

\[
VAN_{E,\text{aggiornato}}
=
VAN_E-\frac{10.000{,}00}{1{,}10}.
\]

\[
VAN_{E,\text{aggiornato}}
\simeq
32.650{,}00-9.090{,}91.
\]

\[
\boxed{
VAN_{E,\text{aggiornato}}
\simeq23.559{,}09
}
\]

\noindent Confrontiamo con zero:

\[
23.559{,}09>0.
\]

\[
\boxed{
\text{E resta accettabile}
}
\]


\subsubsection*{3. Progetto F}

\[
VAN_{F,\text{aggiornato}}
=
VAN_F-\frac{10.000{,}00}{1{,}10}.
\]

\[
VAN_{F,\text{aggiornato}}
\simeq
35.960{,}00-9.090{,}91.
\]

\[
\boxed{
VAN_{F,\text{aggiornato}}
\simeq26.869{,}09
}
\]

\noindent Confrontiamo con zero:

\[
26.869{,}09>0.
\]

\[
\boxed{
\text{F resta accettabile}
}
\]


\subsubsection*{Risposta al punto c}

\begin{center}
    \small
    \begin{tabular}{|c|r|r|c|}
        \hline
        \textbf{Progetto}
        & \textbf{VAN precedente}
        & \textbf{VAN aggiornato}
        & \textbf{Accettabile} \\
        \hline
        E & 32.650,00 & 23.559,09 & Sì \\
        F & 35.960,00 & 26.869,09 & Sì \\
        \hline
    \end{tabular}
\end{center}

\[
\boxed{
\text{Entrambi i progetti: E e F}
}
\]

\noindent Nel punto b il confronto
era con 30.000,00.

\noindent Nel punto c il confronto
è invece con zero:

\[
\boxed{
VAN_{\text{aggiornato}}>0
}
\]

\noindent Entrambi restano accettabili,
anche se il VAN aggiornato
è inferiore a 30.000,00.


\subsection{Punto d:
scelta del paniere attraverso il PI}

\subsubsection*{Che cosa viene richiesto}

\noindent Torniamo ai progetti
del \textbf{punto a}:

\[
B,\qquad E,\qquad F.
\]

\noindent Il budget massimo è:

\[
\boxed{150.000{,}00}
\]

\noindent Seguiamo il metodo del professore:
calcoliamo il PI dei progetti
e confrontiamo i panieri ammissibili.

\noindent Utilizziamo i VAN ottenuti
con CAF pari a 3,79,
senza applicare la spesa del punto c.


\subsubsection*{1. PI del progetto B}

\noindent Applichiamo:

\[
PI_B=\frac{VAN_B}{I_{0,B}}.
\]

\noindent Sostituiamo i dati:

\[
PI_B
\simeq
\frac{14.225{,}00}{90.000{,}00}.
\]

\noindent Espress- oefd % should remove typo
\noindent Esprimendo il rapporto in percentuale:

\[
PI_B
\simeq
\frac{14.225{,}00}{90.000{,}00}\cdot100\%
=
\boxed{15{,}81\%}.
\]

\noindent Per ogni 100 investiti,
B genera un VAN di circa 15,81.


\subsubsection*{2. PI del progetto E}

\[
PI_E=\frac{VAN_E}{I_{0,E}}.
\]

\[
PI_E
\simeq
\frac{32.650{,}00}{100.000{,}00}.
\]

\noindent In percentuale:

\[
PI_E
\simeq
\frac{32.650{,}00}{100.000{,}00}\cdot100\%
=
\boxed{32{,}65\%}.
\]


\subsubsection*{3. PI del progetto F}

\[
PI_F=\frac{VAN_F}{I_{0,F}}.
\]

\[
PI_F
\simeq
\frac{35.960{,}00}{55.000{,}00}.
\]

\noindent In percentuale:

\[
PI_F
\simeq
\frac{35.960{,}00}{55.000{,}00}\cdot100\%
=
\boxed{65{,}38\%}.
\]


\subsubsection*{4. Riepilogo dei PI}

\begin{center}
    \begin{tabular}{|c|r|r|r|}
        \hline
        \textbf{Progetto}
        & \textbf{Investimento}
        & \textbf{VAN}
        & \textbf{PI} \\
        \hline
        B & 90.000,00 & 14.225,00 & 15,81\% \\
        E & 100.000,00 & 32.650,00 & 32,65\% \\
        F & 55.000,00 & 35.960,00 & 65,38\% \\
        \hline
    \end{tabular}
\end{center}

\noindent Il progetto F ha il PI più alto.

\noindent Se scegliamo F,
rimane disponibile un budget di:

\[
150.000{,}00-55.000{,}00
=
95.000{,}00.
\]

\noindent Possiamo aggiungere B,
che richiede 90.000,00,
ma non E,
che richiede 100.000,00.


\subsubsection*{5. Verifichiamo tutte le combinazioni}

\begin{center}
    \small
    \begin{tabular}{|c|r|c|}
        \hline
        \textbf{Paniere}
        & \textbf{Investimento totale}
        & \textbf{Entro il budget} \\
        \hline
        B & 90.000,00 & Sì \\
        E & 100.000,00 & Sì \\
        F & 55.000,00 & Sì \\
        B+E & 190.000,00 & No \\
        B+F & 145.000,00 & Sì \\
        E+F & 155.000,00 & No \\
        B+E+F & 245.000,00 & No \\
        \hline
    \end{tabular}
\end{center}

\noindent I panieri ammissibili sono:

\[
B,\qquad E,\qquad F,\qquad B+F.
\]

\noindent Il confronto richiesto
dal procedimento del professore
è tra E e il paniere B+F.


\subsubsection*{6. Investimento del paniere B+F}

\[
I_{0,B+F}
=
90.000{,}00+55.000{,}00
=
\boxed{145.000{,}00}.
\]

\noindent Il vincolo è rispettato:

\[
145.000{,}00\leq150.000{,}00.
\]

\noindent Il budget inutilizzato è:

\[
150.000{,}00-145.000{,}00
=
\boxed{5.000{,}00}.
\]


\subsubsection*{7. PI del paniere B+F}

\noindent Utilizziamo la media ponderata,
con gli investimenti iniziali come pesi:

\[
PI_{B+F}
=
\frac{
PI_B\cdot90.000{,}00
+
PI_F\cdot55.000{,}00
}{
145.000{,}00
}.
\]

\noindent Sostituendo i PI percentuali
mostrati con due decimali:

\[
PI_{B+F}
\simeq
\frac{
15{,}81\%\cdot90.000{,}00
+
65{,}38\%\cdot55.000{,}00
}{
145.000{,}00
}.
\]

\noindent I due prodotti sono:

\[
15{,}81\%\cdot90.000{,}00
=
14.229{,}00,
\]

\[
65{,}38\%\cdot55.000{,}00
=
35.959{,}00.
\]

\noindent Quindi:

\[
PI_{B+F}
\simeq
\frac{14.229{,}00+35.959{,}00}
{145.000{,}00}.
\]

\[
PI_{B+F}
\simeq
\frac{50.188{,}00}{145.000{,}00}
\simeq
\boxed{34{,}61\%}.
\]

\medskip

\noindent \textbf{Attenzione agli arrotondamenti:}
i PI 15,81\% e 65,38\%
sono già arrotondati.

\noindent Per ritrovare esattamente
il VAN complessivo del paniere,
utilizziamo i rapporti prima
dell'arrotondamento:

\[
PI_{B+F}
=
\frac{
\left(\frac{14.225{,}00}{90.000{,}00}\right)
90.000{,}00
+
\left(\frac{35.960{,}00}{55.000{,}00}\right)
55.000{,}00
}{
145.000{,}00
}.
\]

\noindent Semplificando:

\[
PI_{B+F}
=
\frac{
14.225{,}00+35.960{,}00
}{
145.000{,}00
}.
\]

\[
PI_{B+F}
=
\frac{50.185{,}00}{145.000{,}00}
\simeq
\boxed{34{,}61\%}.
\]

\noindent Il PI percentuale finale
è lo stesso con due decimali.


\subsubsection*{8. Confronto tra E e B+F}

\[
PI_E=32{,}65\%,
\]

\[
PI_{B+F}\simeq34{,}61\%.
\]

\noindent Poiché:

\[
34{,}61\%>32{,}65\%,
\]

\noindent il paniere B+F
presenta un PI maggiore
rispetto al progetto E.


\subsubsection*{9. VAN complessivo del paniere}

\noindent I progetti sono indipendenti,
quindi i loro VAN si sommano:

\[
VAN_{B+F}=VAN_B+VAN_F.
\]

\[
VAN_{B+F}
\simeq
14.225{,}00+35.960{,}00.
\]

\[
\boxed{
VAN_{B+F}\simeq50.185{,}00
}
\]

\noindent Il paniere B+F
ha anche il VAN totale più alto
tra tutte le combinazioni ammissibili:

\[
\begin{aligned}
VAN_B&\simeq14.225{,}00;\\
VAN_E&\simeq32.650{,}00;\\
VAN_F&\simeq35.960{,}00;\\
VAN_{B+F}&\simeq50.185{,}00.
\end{aligned}
\]

\noindent Anche scegliere soltanto F
sarebbe possibile e darebbe un PI più alto,
ma lascerebbe inutilizzata
una parte maggiore del budget.

\noindent Aggiungere B a F
è finanziariamente possibile
e aumenta il VAN totale di 14.225,00.

\noindent Per questo scegliamo B+F:
non stiamo cercando soltanto
il massimo PI percentuale
del singolo progetto.


\subsubsection*{Risposta al punto d}

\[
\boxed{
\text{Paniere scelto: B+F}
}
\]

\[
\boxed{
\begin{aligned}
\text{Investimento totale}&=145.000{,}00;\\
PI_{B+F}&\simeq34{,}61\%;\\
VAN_{B+F}&\simeq50.185{,}00.
\end{aligned}
}
\]


\subsection{Punto e:
recupero dell'investimento entro tre anni}

\subsubsection*{Che cosa viene richiesto}

\noindent Dobbiamo verificare
quali progetti recuperano
l'investimento iniziale
entro i primi tre anni,
utilizzando i flussi attualizzati.

\noindent Il testo richiede
i progetti del punto b:

\[
E,\qquad F.
\]

\noindent Nello svolgimento del professore
compare anche B:
lo riportiamo come confronto,
ma era già escluso dal punto b.

\medskip

\noindent Seguiamo il procedimento
del \textbf{saldo cumulato}:

\begin{enumerate}
    \item partiamo dall'investimento
    con segno negativo;

    \item aggiungiamo il flusso attualizzato
    del primo anno;

    \item aggiungiamo quello del secondo anno;

    \item aggiungiamo quello del terzo anno;

    \item controlliamo quando il saldo
    raggiunge o supera zero.
\end{enumerate}

\noindent Utilizziamo:

\[
\boxed{
S_t=S_{t-1}+\frac{R}{(1{,}10)^t}
}
\]

\noindent I risultati sono mostrati
con due cifre decimali.

\noindent Per evitare errori nella somma,
manteniamo la precisione interna.
Piccole differenze rispetto ai valori
scritti a mano dipendono
dagli arrotondamenti.


\subsubsection*{1. Progetto B:
confronto con la tabella del professore}

\noindent I dati sono:

\[
I_{0,B}=90.000{,}00,
\qquad
R_B=27.500{,}00.
\]

\noindent Partiamo dall'investimento:

\[
S_{B,0}=-90.000{,}00.
\]

\noindent Primo anno:

\[
VA_{B,1}
=
\frac{27.500{,}00}{1{,}10}
=
25.000{,}00.
\]

\[
S_{B,1}
=
-90.000{,}00+25.000{,}00
=
\boxed{-65.000{,}00}.
\]

\noindent Secondo anno:

\[
VA_{B,2}
=
\frac{27.500{,}00}{(1{,}10)^2}
\simeq22.727{,}27.
\]

\[
S_{B,2}
=
S_{B,1}
+
\frac{27.500{,}00}{(1{,}10)^2}
\simeq
\boxed{-42.272{,}73}.
\]

\noindent Terzo anno:

\[
VA_{B,3}
=
\frac{27.500{,}00}{(1{,}10)^3}
\simeq20.661{,}16.
\]

\[
S_{B,3}
=
S_{B,2}
+
\frac{27.500{,}00}{(1{,}10)^3}
\simeq
\boxed{-21.611{,}57}.
\]

\noindent Il saldo dopo tre anni
è ancora negativo.

\[
\boxed{
\text{B non recupera l'investimento entro tre anni}
}
\]


\subsubsection*{2. Progetto E}

\noindent I dati sono:

\[
I_{0,E}=100.000{,}00,
\qquad
R_E=35.000{,}00.
\]

\noindent Partiamo dall'investimento:

\[
S_{E,0}=-100.000{,}00.
\]

\medskip

\noindent \textbf{Primo anno}

\[
VA_{E,1}
=
\frac{35.000{,}00}{1{,}10}
\simeq31.818{,}18.
\]

\[
S_{E,1}
=
-100.000{,}00
+
\frac{35.000{,}00}{1{,}10}
\simeq
\boxed{-68.181{,}82}.
\]

\noindent Il saldo è negativo:
l'investimento non è ancora recuperato.

\medskip

\noindent \textbf{Secondo anno}

\[
VA_{E,2}
=
\frac{35.000{,}00}{(1{,}10)^2}
\simeq28.925{,}62.
\]

\[
S_{E,2}
=
S_{E,1}
+
\frac{35.000{,}00}{(1{,}10)^2}
\simeq
\boxed{-39.256{,}20}.
\]

\noindent Anche dopo due anni
il saldo è negativo.

\medskip

\noindent \textbf{Terzo anno}

\[
VA_{E,3}
=
\frac{35.000{,}00}{(1{,}10)^3}
\simeq26.296{,}02.
\]

\[
S_{E,3}
=
S_{E,2}
+
\frac{35.000{,}00}{(1{,}10)^3}
\simeq
\boxed{-12.960{,}18}.
\]

\noindent Dopo tre anni
mancano ancora circa 12.960,18
per recuperare l'investimento.

\[
\boxed{
\text{E non recupera l'investimento entro tre anni}
}
\]


\subsubsection*{3. Progetto F}

\noindent I dati sono:

\[
I_{0,F}=55.000{,}00,
\qquad
R_F=24.000{,}00.
\]

\noindent Partiamo dall'investimento:

\[
S_{F,0}=-55.000{,}00.
\]

\medskip

\noindent \textbf{Primo anno}

\[
VA_{F,1}
=
\frac{24.000{,}00}{1{,}10}
\simeq21.818{,}18.
\]

\[
S_{F,1}
=
-55.000{,}00
+
\frac{24.000{,}00}{1{,}10}
\simeq
\boxed{-33.181{,}82}.
\]

\noindent Il saldo è ancora negativo.

\medskip

\noindent \textbf{Secondo anno}

\[
VA_{F,2}
=
\frac{24.000{,}00}{(1{,}10)^2}
\simeq19.834{,}71.
\]

\[
S_{F,2}
=
S_{F,1}
+
\frac{24.000{,}00}{(1{,}10)^2}
\simeq
\boxed{-13.347{,}11}.
\]

\noindent Anche dopo due anni
il saldo è negativo:
il recupero non è ancora avvenuto.

\medskip

\noindent \textbf{Terzo anno}

\[
VA_{F,3}
=
\frac{24.000{,}00}{(1{,}10)^3}
\simeq18.031{,}56.
\]

\[
S_{F,3}
=
S_{F,2}
+
\frac{24.000{,}00}{(1{,}10)^3}
\simeq
\boxed{4.684{,}45}.
\]

\noindent Il saldo diventa positivo.

\noindent Con l'incasso del terzo anno,
il progetto ha recuperato
l'investimento iniziale attualizzato
e presenta un margine positivo.

\[
\boxed{
\text{F recupera l'investimento al terzo anno}
}
\]


\subsubsection*{Tabella dei flussi attualizzati}

\begin{center}
    \small
    \begin{tabular}{|l|r|r|r|}
        \hline
        \textbf{Voce}
        & \textbf{B}
        & \textbf{E}
        & \textbf{F} \\
        \hline
        Investimento iniziale
        & -90.000,00 & -100.000,00 & -55.000,00 \\
        \hline
        Flusso annuo
        & 27.500,00 & 35.000,00 & 24.000,00 \\
        \hline
        Flusso attualizzato anno 1
        & 25.000,00 & 31.818,18 & 21.818,18 \\
        \hline
        Flusso attualizzato anno 2
        & 22.727,27 & 28.925,62 & 19.834,71 \\
        \hline
        Flusso attualizzato anno 3
        & 20.661,16 & 26.296,02 & 18.031,56 \\
        \hline
    \end{tabular}
\end{center}


\subsubsection*{Tabella dei saldi cumulati}

\begin{center}
    \small
    \begin{tabular}{|l|r|r|r|}
        \hline
        \textbf{Momento}
        & \textbf{B}
        & \textbf{E}
        & \textbf{F} \\
        \hline
        Tempo iniziale
        & -90.000,00 & -100.000,00 & -55.000,00 \\
        \hline
        Fine anno 1
        & -65.000,00 & -68.181,82 & -33.181,82 \\
        \hline
        Fine anno 2
        & -42.272,73 & -39.256,20 & -13.347,11 \\
        \hline
        Fine anno 3
        & -21.611,57 & -12.960,18 & 4.684,45 \\
        \hline
    \end{tabular}
\end{center}

\noindent \textbf{Come si legge la tabella:}
cerchiamo il primo saldo
uguale o maggiore di zero.

\noindent Soltanto F presenta
un saldo positivo entro il terzo anno.


\subsubsection*{Risposta al punto e}

\[
\boxed{
\text{Si sceglie soltanto F}
}
\]

\noindent E non recupera entro tre anni.

\noindent B, riportato soltanto
per confronto con lo svolgimento,
non recupera entro tre anni
ed era già escluso dal punto b.


\subsection{Riepilogo delle risposte}

\begin{center}
    \begin{tabular}{|c|p{9cm}|}
        \hline
        \textbf{Punto} & \textbf{Risposta} \\
        \hline
        a &
        B, E e F:
        hanno $TIR>10\%$. \\
        \hline
        b &
        E e F:
        hanno $VAN\geq30.000{,}00$,
        utilizzando CAF pari a 3,79. \\
        \hline
        c &
        E e F:
        i VAN aggiornati sono positivi.
        E circa 23.559,09;
        F circa 26.869,09. \\
        \hline
        d &
        B+F:
        investimento 145.000,00;
        $PI\simeq34{,}61\%$;
        $VAN\simeq50.185{,}00$. \\
        \hline
        e &
        Soltanto F:
        il saldo cumulato attualizzato
        diventa positivo al terzo anno,
        raggiungendo circa 4.684,45. \\
        \hline
    \end{tabular}
\end{center}
